\NSPage{013}%
\begin{EESegment}{EE-TGT-P013-001}
The correction cycle now deals with all the linear and quadratic interactions that are still
left, including the curl and cutoff errors in
\NSi{NS-F-P013-001}{\mathscr R(u_B+w,p_B+\pi)}.
\end{EESegment}

\begin{figure}[htbp]
  \centering
  \NSFigurePlaceholder{NS-FIG-5}
  \begin{EESegment}{EE-TGT-P013-002}
  \caption{How the background and the first oscillatory correction are built.
  The averaged pulse flux cancels the divergence of the leading background stress.
  The cycle in Figure~\ref{fig:6} corrects the residual that is left. The flat remainder
  and all curl and cutoff terms stay in the calculation.}
  \label{fig:5}
  \end{EESegment}
\end{figure}

\begin{EESegment}{EE-TGT-P013-003}
\subsection{Correcting the residual and adding up the corrections}\label{sec:3.4}
\end{EESegment}

\begin{EESegment}{EE-TGT-P013-004}
The average wave flux is chosen to cancel the divergence of the leading stress that comes
from \NSi{NS-F-P013-002}{T} (Propositions~\ref{prop:7.5} and~\ref{prop:9.5}). This leaves the
remaining angular harmonics, auxiliary means, and radial integral defects, and each
of these needs its own correction. The corrections always follow the same cycle. After each
operation, the full residual is calculated again, as in Proposition~\ref{prop:9.6}, so any
new terms created by that operation go into the next stage.
\end{EESegment}

\begin{EESegment}{EE-TGT-P013-005}
Start with \NSi{NS-F-P013-003}{(u_B+w,p_B+\pi)}. First reconstruct the pressure, then correct
the angular average of the velocity and the errors in its radial integrals.
Proposition~\ref{prop:9.5} proves that the resulting pair,
\NSi{NS-F-P013-004}{(u^{[0]},p^{[0]})}, has all the bounds needed to begin the correction
cycle. Square brackets count the number of completed cycles. The superscript
\NSi{NS-F-P013-005}{(0)} still marks the leading field
\NSi{NS-F-P013-006}{(u^{(0)},p^{(0)})}.
\end{EESegment}

\begin{EESegment}{EE-TGT-P013-006}
At stage \NSi{NS-F-P013-007}{j}, add a divergence-free velocity increment
\NSi{NS-F-P013-008}{\delta u_j} and a pressure increment
\NSi{NS-F-P013-009}{\delta p_j}:
\NSFormula{NS-F-P013-010}{display}{none}
\[
u^{[j+1]}=u^{[j]}+\delta u_j,
\qquad
p^{[j+1]}=p^{[j]}+\delta p_j.
\]
\end{EESegment}

\begin{EESegment}{EE-TGT-P013-007}
The residual then changes by
\NSFormula{NS-F-P013-011}{display}{none}
\[
\begin{aligned}
\mathscr R(u^{[j+1]},p^{[j+1]})
  &=\mathscr R(u^{[j]},p^{[j]})+\mathcal L_{u^{[j]}}(\delta u_j,\delta p_j)\\
  &\qquad{}+\nabla\mathbin{\cdot}(\delta u_j\otimes\delta u_j).
\end{aligned}
\]
\end{EESegment}

\NSPage{014}%
\begin{EESegment}{EE-TGT-P014-001}
Here \NSi{NS-F-P014-001}{\mathcal L_v} is the linear residual operator defined earlier, with
\NSi{NS-F-P014-002}{v} as its background. So cancelling one chosen source creates
new linear remainders and new quadratic interactions. The four operations in
Proposition~\ref{prop:9.6} handle those terms in the following order.
\end{EESegment}

\begin{EESegment}{EE-TGT-P014-002}
\begin{enumerate}[label=\arabic*.,leftmargin=*]
\item The first operation solves the inhomogeneous wave-amplitude equations for the
supported nonzero angular Fourier modes, the terms in the angular Fourier expansion whose integer
frequency is not zero. The principal operator is the part of the linearized residual that acts on
one new amplitude and its pressure at the current harmonic. It is written explicitly in
\eqref{eq:9.1}, and solving its equation cancels the chosen source. The remaining linear terms and the interactions between the new pulses go
into the next residual (Proposition~\ref{prop:9.1} and Lemma~\ref{lem:9.2}).
\item The second operation builds signed changes in amplitude from the fixed positive
leading amplitudes. For an amplitude change and a leading pulse, their symmetrized cross
covariance adds two doubly averaged products, the leading radial component times the changed
tangential component and the changed radial component times the leading tangential component.
That sum gives the chosen correction to the average stress. Their interactions with earlier
corrections, along with their quadratic self-interaction, remain as higher-order
terms (Proposition~\ref{prop:7.6}, \eqref{eq:9.13}).
\item The third operation corrects the angular average of the residual after its auxiliary
average has been made zero. It does this by inverting the fast auxiliary-time derivative in
\eqref{eq:8.20}. The axial increment comes from the vector potential in \eqref{eq:8.14},
which also gives the radial velocity that goes with it. The reconstruction leaves a remainder
whose Cartesian space-time derivatives go to zero faster than every power of
\NSi{NS-F-P014-003}{q}; that remainder stays in the residual.
\item The fourth operation solves the five radial moment equations in \eqref{eq:8.25}.
Two of them keep the angular-momentum and axial-flux integrals in \eqref{eq:9.10} equal to
zero. The other three cancel the linear contributions to
\NSi{NS-F-P014-004}{(\mathcal P,\mathcal J_\theta,\mathcal J_z)}. These are the errors in the
radial integrals for pressure and tangential momentum defined in \eqref{eq:8.12} and
\eqref{eq:8.15}. Their nonlinear remainders decay faster, as \eqref{eq:8.27} states. The
radial integral constructions also keep the velocity and pressure corrections inside a
compact support.
\end{enumerate}
\end{EESegment}

\begin{EESegment}{EE-TGT-P014-003}
After each operation, the pressure is reconstructed from the updated radial equation, and
its part in the axial equation is included in the new residual. Every product uses the
updated fields, including the curl and cutoff corrections. The background, the leading
amplitudes, and the inverse operators stay fixed for the whole induction.
\end{EESegment}

\begin{EESegment}{EE-TGT-P014-004}
The parameter \NSi{NS-F-P014-005}{\sigma_j} in \eqref{eq:9.8} records how fast the residual
decays. For a Cartesian derivative of order \NSi{NS-F-P014-006}{m}, the bound in
\eqref{eq:9.18} contains the power
\NSi{NS-F-P014-007}{q^{h\sigma_j-K_m}}. Here
\NSi{NS-F-P014-008}{K_m} does not depend on
\NSi{NS-F-P014-009}{j}. The same bound also contains logarithmic factors and the remainder
that is being kept separate. The constants and logarithmic powers may depend on
\NSi{NS-F-P014-010}{j,m}. At the start, and after each correction cycle, the values are
\NSFormula{NS-F-P014-011}{display}{none}
\[
\sigma_0=\frac15,
\qquad
\sigma_{j+1}=\sigma_j+\frac1{10},
\qquad
\sigma_j=\frac15+\frac j{10}\longrightarrow\infty.
\]
The cumulative velocity bounds in \eqref{eq:9.9}, the support conditions, and the two
preserved integrals in \eqref{eq:9.10} all continue to hold, so the cycle can be run again.
Lemma~\ref{lem:9.7} gives one common domain on which every finite stage is defined.
\end{EESegment}

\begin{EESegment}{EE-TGT-P014-005}
To add up all the corrections, first put cutoffs on their vector potentials and pressures,
then take curls. Put the same cutoffs on the direct axisymmetric azimuthal increments. Each
cutoff is equal to one close enough to \NSi{NS-F-P014-012}{q=0}, and its support gets smaller
as the stage increases. The bounds in \eqref{eq:9.17} lose a power of
\NSi{NS-F-P014-013}{q} that depends on the order of the derivative but not on the stage.
Because of these bounds, the cutoffs can be chosen so that every differentiated tail obeys
\eqref{eq:5.35}. At any point with \NSi{NS-F-P014-014}{q>0}, only finitely many terms in the
sum are nonzero. The sum gives smooth fields
\NSi{NS-F-P014-015}{(u_{\mathrm{loc}},p_{\mathrm{loc}})} for
\NSi{NS-F-P014-016}{t<1}. Compare their residual with the residual at any one fixed finite
stage to get \eqref{eq:9.20}. The leading growth of the velocity is not changed
(Proposition~\ref{prop:9.9}).
\end{EESegment}

\begin{EESegment}{EE-TGT-P014-006}
This sum gives the needed local growth of the velocity. Passing from the local construction
to the whole space also needs a vector-potential representation and regularity through time
one everywhere away from the origin. Proposition~\ref{prop:9.9} proves both facts, along with
the flatness of the residual. In the theorem, \NSi{NS-F-P014-017}{(u,p)} means the local pair
\NSi{NS-F-P014-018}{(u_{\mathrm{loc}},p_{\mathrm{loc}})}.
\end{EESegment}

\NSPage{015}%
\begin{figure}[htbp]
  \centering
  \NSFigurePlaceholder{NS-FIG-6}
  \begin{EESegment}{EE-TGT-P015-001}
  \caption{Correcting the residual and finishing the construction, continuing
  Figure~\ref{fig:5}. Each turn of the correction cycle raises the residual exponent
  \NSi{NS-F-P015-001}{\sigma_j}, while every linear and nonlinear term stays in the
  calculation. The flatness estimate holds uniformly for
  \NSi{NS-F-P015-002}{0\leq X\leq X_1} as
  \NSi{NS-F-P015-003}{q\downarrow0}, for each finite
  \NSi{NS-F-P015-004}{X_1}. The velocity stays incompressible through the whole
  construction.}
  \label{fig:6}
  \end{EESegment}
\end{figure}

\begin{EESegment}{EE-TGT-P015-002}
The notation \NSi{NS-F-P015-005}{\partial_x^\alpha\partial_t^b} means Cartesian derivatives
in space and time. Here \NSi{NS-F-P015-006}{\alpha} is a multi-index for the spatial
derivatives, and \NSi{NS-F-P015-007}{b\geq0} is an integer giving the number of time
derivatives.
\end{EESegment}

\begin{theorem}\label{thm:3.1}
\begin{EESegment}{EE-TGT-P015-003}
There are constants \NSi{NS-F-P015-008}{0<h<1/100},
\NSi{NS-F-P015-009}{q_*>0}, \NSi{NS-F-P015-010}{0<X_a<X_{\mathrm{ext}}<\infty}, and smooth
fields on
\NSFormula{NS-F-P015-011}{display}{none}
\[
\Omega_* = \{(x,t):\tau>0,\ q<q_*\}
\]
with the following properties.
\end{EESegment}
\begin{enumerate}[label=\textup{(\roman*)},leftmargin=*]
\item
\begin{EESegment}{EE-TGT-P015-004}
There is a vector potential \NSi{NS-F-P015-012}{\mathcal A} and an axisymmetric scalar
\NSi{NS-F-P015-013}{\mathcal B} such that
\NSFormula{NS-F-P015-014}{display}{3.3}
\begin{equation}
u=\operatorname{curl}\mathcal A+\mathcal B e_\theta.
\tag{3.3}\label{eq:3.3}
\end{equation}
Both \NSi{NS-F-P015-015}{\mathcal A} and \NSi{NS-F-P015-016}{\mathcal B e_\theta} are smooth
in Cartesian coordinates at the axis. The field is exactly divergence-free.
\end{EESegment}

\item
\begin{EESegment}{EE-TGT-P015-005}
Every Cartesian space-time derivative of
\NSi{NS-F-P015-017}{\mathcal A,\mathcal B e_\theta,p} stays uniformly bounded on every
compact spatial set where \NSi{NS-F-P015-018}{c\leq q\leq c'<q_*}, for each
\NSi{NS-F-P015-019}{0<c<c'}, all the way up to the one-sided boundary
\NSi{NS-F-P015-020}{\tau=0}. The limits of these derivatives
agree with one another when another spatial or time derivative is taken.
\end{EESegment}

\item
\begin{EESegment}{EE-TGT-P015-006}
For every \NSi{NS-F-P015-021}{\alpha,b,N} and every finite
\NSi{NS-F-P015-022}{X_1},
\NSFormula{NS-F-P015-023}{display}{3.4}
\begin{equation}
\lvert\partial_x^\alpha\partial_t^b\mathscr R(u,p)\rvert
\leq C_{\alpha,b,N,X_1}q^N
\qquad (0\leq X\leq X_1,\ q\downarrow0).
\tag{3.4}\label{eq:3.4}
\end{equation}
For \NSi{NS-F-P015-024}{X\geq X_{\mathrm{ext}}}, the residual is exactly zero, and
\NSFormula{NS-F-P015-025}{display}{3.5}
\begin{equation}
\mathcal A=0,
\qquad \mathcal B=K(r,\tau),
\qquad p=-\int_r^\infty\frac{K(\rho,\tau)^2}{\rho}\,d\rho,
\qquad K(r,\tau)=r^{-1-2h}\mathcal H_{\mathrm{ext}}(\tau/r^2).
\tag{3.5}\label{eq:3.5}
\end{equation}
Here \NSi{NS-F-P015-026}{\mathcal H_{\mathrm{ext}}} is smooth on
\NSi{NS-F-P015-027}{[0,\infty)}, and each one of its fixed derivatives is bounded. The
radial swirl heat equation obeyed by \NSi{NS-F-P015-028}{K} is
\NSFormula{NS-F-P015-029}{display}{none}
\[
-\partial_\tau K=\partial_r^2K+r^{-1}\partial_rK-r^{-2}K.
\]
\end{EESegment}

\NSPage{016}%
\begin{EESegment}{EE-TGT-P016-001}
In this statement, the second input of \NSi{NS-F-P016-001}{K} is
\NSi{NS-F-P016-002}{\tau}. Later, the notation \NSi{NS-F-P016-003}{K(r,t)} means the same
field evaluated at \NSi{NS-F-P016-004}{\tau=1-t}.
\end{EESegment}

\item
\begin{EESegment}{EE-TGT-P016-002}
There are fixed values \NSi{NS-F-P016-005}{X_{\mathrm{in}}\in(0,X_a)} and
\NSi{NS-F-P016-006}{e_0>0} for which
\NSFormula{NS-F-P016-007}{display}{3.6}
\begin{equation}
u_\theta(\sqrt{2X_{\mathrm{in}}\tau},0,0,1-\tau)
=\tau^{-A}\bigl(e_0+O(\tau^{2h})\bigr)
\qquad(\tau\downarrow0).
\tag{3.6}\label{eq:3.6}
\end{equation}
The four entries on the left give the cylindrical position and the time.
\end{EESegment}
\end{enumerate}
\end{theorem}

\begin{EESegment}{EE-TGT-P016-003}
\subsection{Localizing the construction and finishing the proof}\label{sec:3.5}
\end{EESegment}

\begin{EESegment}{EE-TGT-P016-004}
The last step uses the decomposition in \eqref{eq:3.3}. Multiply the vector potential
\NSi{NS-F-P016-008}{\mathcal A}, the azimuthal coefficient
\NSi{NS-F-P016-009}{\mathcal B}, and the pressure
\NSi{NS-F-P016-010}{p_{\mathrm{loc}}} by \NSi{NS-F-P016-011}{c=\chi_x\chi_t}. The spatial
cutoff \NSi{NS-F-P016-012}{\chi_x} is smooth and axisymmetric and has compact support
\NSi{NS-F-P016-013}{\mathcal K}. The time cutoff \NSi{NS-F-P016-014}{\chi_t} is smooth and
satisfies \NSi{NS-F-P016-015}{\chi_t=0} on \NSi{NS-F-P016-016}{[0,1-\tau_0]} for some
\NSi{NS-F-P016-017}{0<\tau_0<1/2}. Both cutoffs are equal to one near
\NSi{NS-F-P016-018}{(0,1)}. Since \NSi{NS-F-P016-019}{q\asymp\tau+|z|^{1/D}}, their
supports can be chosen so that \NSi{NS-F-P016-020}{q<q_*/2} wherever
\NSi{NS-F-P016-021}{c\neq0} and \NSi{NS-F-P016-022}{t<1}. After multiplying, take the curl
and extend the fields by zero outside the support of the cutoffs:
\NSFormula{NS-F-P016-023}{display}{none}
\[
u=\operatorname{curl}(c\mathcal A)+c\mathcal B e_\theta,
\qquad p=cp_{\mathrm{loc}},
\qquad
u=cu_{\mathrm{loc}}+\nabla c\times\mathcal A.
\]
The curl term and the axisymmetric azimuthal term are each divergence-free. The chosen
representatives are smooth in Cartesian coordinates, and their support stays away from
\NSi{NS-F-P016-024}{q=q_*}, so extending them by zero is smooth. The fields
\NSi{NS-F-P016-025}{u,p} have the fixed spatial support
\NSi{NS-F-P016-026}{\mathcal K}, vanish throughout the stated initial time interval, and
agree with the local fields near \NSi{NS-F-P016-027}{(0,1)}
(Proposition~\ref{prop:10.1}).
\end{EESegment}

\begin{EESegment}{EE-TGT-P016-005}
For \NSi{NS-F-P016-028}{t<1}, set \NSi{NS-F-P016-029}{f=\mathscr R(u,p)}. When
\NSi{NS-F-P016-030}{f-c\mathscr R(u_{\mathrm{loc}},p_{\mathrm{loc}})} is expanded,
differentiating the cutoffs gives terms such as
\NSi{NS-F-P016-031}{(\partial_t c-\Delta c)u_{\mathrm{loc}}} and
\NSi{NS-F-P016-032}{p_{\mathrm{loc}}\nabla c}. The curl correction
\NSi{NS-F-P016-033}{\nabla c\times\mathcal A} adds its derivatives and its advection
products. The change in self-advection also contains
\NSi{NS-F-P016-034}{(c^2-c)(u_{\mathrm{loc}}\mathbin{\cdot}\nabla)u_{\mathrm{loc}}}. All
these extra terms live in the parts where the cutoffs are changing, away from a neighborhood
of \NSi{NS-F-P016-035}{(0,1)}.
\end{EESegment}

\begin{EESegment}{EE-TGT-P016-006}
Near \NSi{NS-F-P016-036}{(0,1)}, the cutoffs are equal to one, so
\NSi{NS-F-P016-037}{f} is the local residual. On
\NSi{NS-F-P016-038}{X\leq X_{\mathrm{ext}}}, its flatness estimate \eqref{eq:3.4}, together
with the exactly zero residual for larger \NSi{NS-F-P016-039}{X}, gives the bound
\eqref{eq:10.9} for every Cartesian space-time
derivative. In the cutoff transition regions near time one,
Theorem~\ref{thm:3.1}\textup{(ii)} bounds every derivative of the actual potentials and
pressure because \NSi{NS-F-P016-040}{q} stays positive there. In the part that remains,
\NSi{NS-F-P016-041}{r} stays positive while \NSi{NS-F-P016-042}{q\to0}. There
\NSi{NS-F-P016-043}{\mathcal A=0}, and the heat bounds in \eqref{eq:10.8}, together with the
normalized pressure integral in \eqref{eq:3.5}, give the limits of the derivatives. These
estimates, the uniform flatness bound near the origin, and the fixed support
\NSi{NS-F-P016-044}{\mathcal K} give compatible limits, uniformly on
\NSi{NS-F-P016-045}{\R^3}, for every derivative of \NSi{NS-F-P016-046}{f}
(Lemma~\ref{lem:10.2}). Lemma~\ref{lem:10.3} realizes those limits from
\NSi{NS-F-P016-047}{t>1} with a force supported in
\NSi{NS-F-P016-048}{\mathcal K\times[0,2]}:
\NSFormula{NS-F-P016-049}{display}{none}
\[
f\in C_c^\infty(\R^3\times(0,\infty);\R^3),
\qquad
\mathscr R(u,p)=f
\qquad(0\leq t<1).
\]
\end{EESegment}

\begin{EESegment}{EE-TGT-P016-007}
The growth path in Theorem~\ref{thm:3.1}\textup{(iv)} eventually stays in the region where
the cutoffs are equal to one, so the localized velocity keeps the asymptotic formula
\eqref{eq:3.6} and becomes unbounded as \NSi{NS-F-P016-050}{t\uparrow1}.
\end{EESegment}

\begin{EESegment}{EE-TGT-P016-008}
The concentration scales still allow the energy to stay bounded. The core
\NSi{NS-F-P016-051}{\mathcal C_\tau} has volume of order
\NSi{NS-F-P016-052}{\tau^{3/2-h}}. Its leading kinetic energy
\NSi{NS-F-P016-053}{E_{\mathrm{core}}}, and the integral of the squared radial derivatives
of its leading velocity, \NSi{NS-F-P016-054}{D_{\mathrm{core}}}, scale as
\NSFormula{NS-F-P016-055}{display}{none}
\[
E_{\mathrm{core}}\asymp\tau^{3/2-h}\tau^{-1-2h}=\tau^{1/2-3h},
\qquad
D_{\mathrm{core}}\asymp\tau^{3/2-h}\tau^{-2-2h}=\tau^{-1/2-3h}.
\]
In each product, the first factor is the volume of the region, and the second is either the
squared velocity or the squared radial derivative. Since \NSi{NS-F-P016-056}{h<1/6},
\NSFormula{NS-F-P016-057}{display}{none}
\[
\tau^{1/2-3h}\longrightarrow0,
\qquad
\int_0^{\tau_0}\tau^{-1/2-3h}\,d\tau
=\frac{\tau_0^{1/2-3h}}{1/2-3h}<\infty
\qquad(\tau_0>0).
\]
\end{EESegment}

\NSPage{017}%
\begin{EESegment}{EE-TGT-P017-001}
Lemma~\ref{lem:10.4} proves the global energy bound and the bound on integrated dissipation
for the localized fields directly from their equation.
\end{EESegment}

\begin{EESegment}{EE-TGT-P017-002}
Lemma~\ref{lem:10.5} now applies to any smooth solution with the same force and zero initial
data whose kinetic energy stays uniformly bounded. That solution must agree with
\NSi{NS-F-P017-001}{u} on every \NSi{NS-F-P017-002}{[0,T]} with
\NSi{NS-F-P017-003}{T<1}. If a global smooth solution existed, it would agree with the
constructed velocity for \NSi{NS-F-P017-004}{0\leq t<1}. Along the growth path, however,
the constructed velocity is not bounded on any compact neighborhood of
\NSi{NS-F-P017-005}{(0,1)}, while a global smooth solution would be. This contradiction
proves Theorem~\ref{thm:1.1}.
\end{EESegment}

\begin{EESegment}{EE-TGT-P017-003}
\subsection{Notation}\label{sec:3.6}
Every choice of profile is fixed before \NSi{NS-F-P017-006}{q\downarrow0}. Constants may
depend on those choices and on the stated order of differentiation, but not on the
concentration scale. The notation \NSi{NS-F-P017-007}{O(q^a)} means that the quantity is
bounded by \NSi{NS-F-P017-008}{Cq^a} on the stated domain. Any uniformity in other
parameters is stated when it is needed. The background expansion order
\NSi{NS-F-P017-009}{n}, the correction stage \NSi{NS-F-P017-010}{j}, the dyadic band index
\NSi{NS-F-P017-011}{\ell}, and the Fourier harmonic index are four different indices. The dyadic
band index labels successive scales, each half the size of the preceding one, and records the scale
to which the concentration parameter is comparable.
\end{EESegment}

\begin{EESegment}{EE-TGT-P017-004}
\paragraph{Averages.}
Let \NSi{NS-F-P017-012}{F(r,\theta,z,t,Y)} be a scalar field on the extended domain. Its
auxiliary variable \NSi{NS-F-P017-013}{Y\in\T^2=\R^2/\Z^2} is independent of the physical
coordinates. The angular average and the auxiliary average are
\NSFormula{NS-F-P017-014}{display}{3.7}
\begin{equation}
\begin{aligned}
\langle F\rangle_\theta(r,z,t,Y)
  &:=\frac1{2\pi}\int_0^{2\pi}F(r,\vartheta,z,t,Y)\,d\vartheta,\\
\langle F\rangle_Y(r,\theta,z,t)
  &:=\int_{\T^2}F(r,\theta,z,t,Y')\,dY',
  \qquad \int_{\T^2}1\,dY'=1.
\end{aligned}
\tag{3.7}\label{eq:3.7}
\end{equation}
The angular average varies the angle while keeping the independent auxiliary coordinate
fixed. The auxiliary average varies that coordinate while keeping all the physical
coordinates fixed. For vector and tensor fields, both averages act on the cylindrical
components. Applying both of them, in either order, gives the fast average used for the wave
fields:
\NSFormula{NS-F-P017-015}{display}{3.8}
\begin{equation}
\langle F\rangle
:=\langle\langle F\rangle_\theta\rangle_Y
=\langle\langle F\rangle_Y\rangle_\theta
=\frac1{2\pi}\int_{\T^2}\int_0^{2\pi}F(r,\vartheta,z,t,Y)\,d\vartheta\,dY.
\tag{3.8}\label{eq:3.8}
\end{equation}
These averages are taken before the field is restricted to the physical phase map in
\eqref{eq:6.3}. They keep \NSi{NS-F-P017-016}{(r,z,t)} fixed, or, in any one normalized
chart, keep \NSi{NS-F-P017-017}{(R,Z,T)} fixed. By \eqref{eq:6.19}, the normalized Haar
averages agree on the absolute auxiliary torus, each band torus, and the common auxiliary
torus. For the wave construction, this two-step average is the whole fast average, including
its normalization. If a coefficient is already an angular mean, its part with auxiliary
average zero is
\NSFormula{NS-F-P017-018}{display}{3.9}
\begin{equation}
f^\circ:=f-\langle f\rangle_Y,
\qquad
\langle f^\circ\rangle_Y=0.
\tag{3.9}\label{eq:3.9}
\end{equation}
\end{EESegment}

\begin{EESegment}{EE-TGT-P017-005}
The profile construction also has a different auxiliary loop angle,
\NSi{NS-F-P017-019}{\theta'}, with period \NSi{NS-F-P017-020}{2\pi}. Its average is
\NSi{NS-F-P017-021}{\langle f\rangle_{\theta'}=(2\pi)^{-1}\int_0^{2\pi}f(\theta')\,d\theta'},
as in Lemmas~\ref{lem:C.1} and~\ref{lem:4.11}. In that local profile notation, brackets with
no subscript mean this loop average. The radial average is
\NSFormula{NS-F-P017-022}{display}{3.10}
\begin{equation}
\mathcal A_X(f)(X,\eta)=\frac1X\int_0^X f(X',\eta)\,dX'.
\tag{3.10}\label{eq:3.10}
\end{equation}
The auxiliary torus supplies the parameters for the oscillations. The torus in
Corollary~\ref{cor:10.6} is a
different object, the spatial domain of the periodic solution.
\end{EESegment}

\begin{EESegment}{EE-TGT-P017-006}
\paragraph{Normalized operators and coefficient classes.}
Inside a dyadic chart, \NSi{NS-F-P017-023}{Q=2^{-\ell}} is fixed,
\NSi{NS-F-P017-024}{q\asymp Q}, and
\NSFormula{NS-F-P017-025}{display}{none}
\[
(R,Z,T)=(r/\sqrt Q,z/Q^D,\tau/Q),
\qquad
\varepsilon=Q^h,
\qquad
S_*=\ell^2.
\]
The star on a spatial operator includes the factor \NSi{NS-F-P017-026}{\sqrt Q} that changes
physical spatial derivatives into chart units. A physical derivative also follows the
auxiliary phase map through the chain rule. With the radial operator
\NSi{NS-F-P017-027}{\mathfrak r} from \eqref{eq:6.4}, set
\NSi{NS-F-P017-028}{D_r=\sqrt Q\,\mathfrak r},
\NSi{NS-F-P017-029}{D_z=\sqrt Q\,\partial_z=\varepsilon\partial_Z}, and
\end{EESegment}

\NSPage{018}%
\begin{EESegment}{EE-TGT-P018-001}
\NSi{NS-F-P018-001}{D_\theta=R^{-1}\partial_\theta}, as in \eqref{eq:6.6}. For a scalar
\NSi{NS-F-P018-002}{f} and a vector
\NSi{NS-F-P018-003}{a=(a_r,a_\theta,a_z)} written in cylindrical components, the normalized
gradient, divergence, and curl are
\NSFormula{NS-F-P018-004}{display}{3.11}
\begin{equation}
\begin{aligned}
\nabla_*f&=(D_rf,D_\theta f,D_zf),\\
\operatorname{div}_*a&=(D_r+R^{-1})a_r+D_\theta a_\theta+D_za_z,\\
\operatorname{curl}_*a&=(D_\theta a_z-D_za_\theta,\ D_za_r-D_ra_z,\ (D_r+R^{-1})a_\theta-D_\theta a_r).
\end{aligned}
\tag{3.11}\label{eq:3.11}
\end{equation}
The separate factors that normalize the fields themselves are given in \eqref{eq:8.11}.
The physical momentum residual is \NSi{NS-F-P018-005}{\mathscr R} from \eqref{eq:3.1}.
\end{EESegment}

\begin{EESegment}{EE-TGT-P018-002}
The classes \NSi{NS-F-P018-006}{\mathsf W_\alpha,\mathsf M_\alpha,\mathsf S_\alpha} keep
track, in normalized units, of wave amplitudes, angular mean coefficients, and radial
moments. A coefficient derivative differentiates the amplitude after its oscillatory
exponential has been factored out, while keeping the band, label, and harmonic fixed. At
each fixed correction stage and each fixed coefficient-derivative order,
\NSi{NS-F-P018-007}{\partial^I} means ordinary differentiation in
\NSi{NS-F-P018-008}{(R,Z,T,H)}. Here \NSi{NS-F-P018-009}{H} is the common auxiliary
coordinate from \eqref{eq:6.17}. For moments, only \NSi{NS-F-P018-010}{(Z,T)} are
differentiated. Their bounds have the following forms. The first two bounds hold on
\NSi{NS-F-P018-011}{X_a<X<X_b}, and the wave bound is also restricted to its labelled
rectangle:
\NSFormula{NS-F-P018-012}{display}{none}
\[
\begin{aligned}
f\in \mathsf M_\alpha
&\quad\Longrightarrow\quad
|\partial^If|\leq C_{j,I}\varepsilon^\alpha S_*^{b_{j,I}}\zeta\delta^{-d_{j,I}},\\
a\in \mathsf W_\alpha
&\quad\Longrightarrow\quad
|\partial^Ia|\leq C_{j,I}\varepsilon^\alpha S_*^{b_{j,I}}\sqrt\zeta\,\delta^{-d_{j,I}}P_v
\qquad(0\leq v\leq L_s),\\
F\in \mathsf S_\alpha
&\quad\Longrightarrow\quad
|\partial_{Z,T}^IF|\leq C_{j,I}\varepsilon^\alpha S_*^{b_{j,I}}.
\end{aligned}
\]
Here \NSi{NS-F-P018-013}{\zeta,\delta} are the weights near the two radial edges from
\eqref{eq:6.23}, and \NSi{NS-F-P018-014}{P_v} is the pulse envelope from
\eqref{eq:7.12}, used as a pointwise weight. The constants
\NSi{NS-F-P018-015}{C_{j,I}>0} and \NSi{NS-F-P018-016}{b_{j,I},d_{j,I}\geq0} can be
different in the three rows. They may depend on the fixed stage, the derivative order, and
the profile choices, but they are uniform over the band, label, copy of the rectangle, and
point. Definitions~\ref{def:6.4} and~\ref{def:6.5} give the full definitions, including the
support, compatibility, and smooth-extension conditions. An angular mean coefficient can
still depend on the auxiliary torus.
\end{EESegment}

\begin{EESegment}{EE-TGT-P018-003}
\begin{definition}[Regularity at the axis]\label{def:3.2}
A leading field is regular at the axis when its velocity and pressure extend smoothly in
Cartesian coordinates across \NSi{NS-F-P018-017}{r=0}, for
\NSi{NS-F-P018-018}{t<1}. For scalar auxiliary profiles, ``regular at the axis'' means that
the profile extends smoothly in \NSi{NS-F-P018-019}{(X,\eta)} to
\NSi{NS-F-P018-020}{X=0}.
\end{definition}
\end{EESegment}

\begin{EESegment}{EE-TGT-P018-004}
\begin{definition}[Small parameters and the order in which they are chosen]\label{def:3.3}
For \NSi{NS-F-P018-021}{\alpha\geq0} and \NSi{NS-F-P018-022}{\beta>0}, the notation
\NSi{NS-F-P018-023}{\alpha\ll\beta} means that
\NSi{NS-F-P018-024}{\alpha/\beta} must be small enough, where how small it must be can
depend on all the data already fixed. In a chain of constants, choose them from right to
left. A small reciprocal means that the corresponding parameter is large. When an
expression depends on profile variables, the stated smallness holds uniformly on its stated
domain and follows from those choices of parameters.
\end{definition}
\end{EESegment}

\begin{EESegment}{EE-TGT-P018-005}
The next table lists the main symbols in the construction and points to their definitions.
It puts the different background fields in one group. When the same letter is used more
than once, the table tells the uses apart by their role or by the section where they appear.
An order attached to a coefficient class refers to its normalized representative, with the
weights and derivative bounds given in the definition of that class.
\end{EESegment}
